\begin{proof}
\textbf{(1) Set the conjugate weights.}  
Define
\[
\alpha:=\frac{1}{p},\qquad \beta:=\frac{1}{q}.
\]
Since $p,q>1$, we have $\alpha,\beta\in(0,1)$ and
\[
\alpha+\beta=\frac{1}{p}+\frac{1}{q}=1 .
\]

\medskip

\textbf{(2) Degenerate cases.}  
If $a=0$ or $b=0$, then $ab=0$ while $\frac{a^{p}}{p}+\frac{b^{q}}{q}\ge0$, so the inequality holds with equality. Hence we may assume $a>0$ and $b>0$ in what follows.

\medskip

\textbf{(3) Weighted AM–GM.}  
For any non‑negative numbers $X,Y$ and weights $\alpha,\beta\ge0$ with $\alpha+\beta=1$, the weighted arithmetic–geometric mean inequality states
\[
\alpha X+\beta Y\;\ge\;X^{\alpha}Y^{\beta}.
\tag{3}
\]
Equality in \eqref{3} occurs iff $X=Y$.

\medskip

\textbf{(4) Apply \eqref{3} with $X=a^{p}$, $Y=b^{q}$.}  
Since $a,b\ge0$, also $X,Y\ge0$. Substituting gives
\[
\alpha a^{p}+\beta b^{q}\;\ge\;(a^{p})^{\alpha}(b^{q})^{\beta}.
\tag{4}
\]

\medskip

\textbf{(5) Simplify the right‑hand side of \eqref{4}.}  
Using $\alpha=1/p$ and $\beta=1/q$,
\[
(a^{p})^{\alpha}(b^{q})^{\beta}
= (a^{p})^{1/p}\,(b^{q})^{1/q}
= a\,b ,
\]
because the bases are non‑negative and we take the principal roots.  Inserting this into \eqref{4} yields
\[
\frac{a^{p}}{p}+\frac{b^{q}}{q}\;\ge\;ab .
\tag{5}
\]

\medskip

\textbf{(6) Conclusion.}  
Inequality \eqref{5} is precisely Young’s inequality. Together with the degenerate case handled in step~2, the statement holds for all non‑negative $a,b$.

\medskip

\textbf{(7) Equality case.}  
Equality in the weighted AM–GM occurs exactly when $X=Y$, i.e. when $a^{p}=b^{q}$. Together with the trivial possibilities $a=0$ or $b=0$, the equality cases are
\[
a=0\quad\text{or}\quad b=0\quad\text{or}\quad a^{p}=b^{q}>0 .
\]
\end{proof}